\documentclass[11pt,letterpaper,oneside]{scrreprt}

\usepackage[margin=1in]{geometry}
\usepackage{iftex}
\ifPDFTeX
  \usepackage[T1]{fontenc}
\else
  \usepackage{fontspec}
  \IfFontExistsTF{Cambria}
    {\setmainfont{Cambria}}
    {\setmainfont{Latin Modern Roman}}
  \IfFontExistsTF{Calibri}
    {\setsansfont{Calibri}}
    {\setsansfont{Latin Modern Sans}}
\fi
\usepackage[unicode,bookmarks=false,linktoc=none]{hyperref}

\hypersetup{
  pdftitle={Dragon Quest V: Comprehensive Guide},
  pdfsubject={Dragon Quest V for Super Famicom / SNES},
  hidelinks
}

\setcounter{tocdepth}{2}
\setcounter{secnumdepth}{2}
\newcommand{\DiscardPageNumber}[1]{}
\DeclareTOCStyleEntries[
  pagenumberbox=\DiscardPageNumber,
  linefill=\hfill
]{tocline}{part,chapter,section,subsection}
\DeclareTOCStyleEntry[indent=0pt,numwidth=2.4em]{tocline}{chapter}
\DeclareTOCStyleEntry[indent=2.4em,numwidth=3.2em]{tocline}{section}
\DeclareTOCStyleEntry[indent=5.6em,numwidth=4.2em]{tocline}{subsection}
\renewcommand{\contentsname}{Table of Contents}

\newcommand{\PlanPart}[1]{%
  \refstepcounter{part}%
  \addcontentsline{toc}{part}{Part \Roman{part}: #1}%
}
\newcommand{\PlanChapter}[2]{%
  \refstepcounter{chapter}%
  \label{#1}%
  \addcontentsline{toc}{chapter}{\protect\numberline{\thechapter}#2}%
}
\newcommand{\PlanSection}[1]{%
  \refstepcounter{section}%
  \addcontentsline{toc}{section}{\protect\numberline{\thesection}#1}%
}
\newcommand{\PlanSubsection}[1]{%
  \refstepcounter{subsection}%
  \addcontentsline{toc}{subsection}{\protect\numberline{\thesubsection}#1}%
}

\begin{document}
\pagestyle{plain}
\begin{titlepage}
\centering
\vspace*{1cm}
{\sffamily\bfseries\Huge Dragon Quest V\par}
\bigskip
{\sffamily\Large Comprehensive Guide\par}
\bigskip
{\large Tenkuu no Hanayome\par}
{\large Super Famicom / SNES\par}
\vfill
\begin{minipage}{\textwidth}
\emph{Dragon Quest V: Tenkuu no Hanayome}\par
Enix, 1992\par
Super Famicom / SNES
\end{minipage}
\vfill
\end{titlepage}

\PlanPart{Foundations and Game Systems}
\PlanChapter{conventions}{Game Version, Names, and Guide Conventions}
\PlanSection{Game Version}
\PlanSection{Names}

\PlanChapter{exploration}{Controls, Exploration, and Progression}
\PlanSection{Controller inputs, field commands, and menu navigation}
\PlanSection{Data and status screens: attributes, experience, and equipment}
\PlanSection{Dialogue, searching, drawers, pots, chests, and hidden objects}
\PlanSection{Day and night: access, shops, events, and changing time}
\PlanSection{Doors, keys, moving obstacles, floor hazards, and puzzles}
\PlanSection{Inventory ownership, item transfer, capacity, and equipment use}
\PlanSection{Inns, churches, saving, revival, and status recovery}
\PlanSection{Banks, gold storage, shopping, and selling}
\PlanSection{Travel progression: walking, wagon, ship, carpet, castle, and dragon}
\PlanSection{Return destinations, Outside, maps, and safe backtracking}

\PlanChapter{combat}{Combat, Tactics, and Party Management}
\PlanSection{Battle commands, target groups, turn order, and escape}
\PlanSection{Physical attacks, critical hits, multi-target weapons, and damage}
\PlanSection{Spells, special abilities, MP, and reusable battle items}
\PlanSection{Status ailments, buffs, debuffs, resistance, and spell reflection}
\PlanSection{AI tactic modes and when to use direct commands}
\PlanSection{Intelligence, obedience, and monster command restrictions}
\PlanSection{The three-member active party and wagon reserves}
\PlanSection{Changing members, wagon access, experience, and dungeon restrictions}
\PlanSection{Healing economy, resource conservation, and boss preparation}
\PlanSection{Enemy summoning, metal enemies, defeat, and recovery}

\PlanPart{Complete Three-Generation Walkthrough}
\PlanChapter{walk-childhood}{First Generation: Childhood}
\PlanSection{Birth scene, the opening ship, and Vista Port}
\PlanSection{Santa Rosa: Sancho, Bianca, supplies, and the missing herbalist}
\PlanSection{Santa Rosa Cave: rescue route, treasure, and early training}
\PlanSection{Alcapa / Alkapa: Bianca's home and the captive panther}
\PlanSection{Lenoire Castle: entry, separated companions, darkness, and the torch}
\PlanSubsection{Ghost-boss preparation and battle strategy}
\PlanSubsection{Golden Orb, the panther's names, and Bianca's departure}
\PlanSection{Santa Rosa revisited: Bella and the route to the faerie realm}
\PlanSection{Faerie Village and the Dwarves Cave: Thieves' Key knowledge}
\PlanSection{Ice Mansion: sliding-floor route, Zaile, and the final boss}
\PlanSection{Flute of Spring, returning home, and childhood loose ends}
\PlanSection{Reinhart Castle: Henry's disappearance and the kidnapping}
\PlanSection{Thieves Cave: Henry's rescue, Gema, and the scripted transition}
\PlanSection{First-generation treasure and pre-transition checklist}

\PlanChapter{walk-escape}{Second Generation: Escape and Reinhart}
\PlanSection{Captivity, Maria, Henry, and escape from the Great Temple}
\PlanSection{The Nameless Nunnery: recovery and the first travel objective}
\PlanSection{Oracleberry / Oracle Valley: shops, wagon, map, and monster storage}
\PlanSubsection{First monster recruits and party-building opportunities}
\PlanSubsection{Casino introduction and optional early prize planning}
\PlanSection{Santa Rosa Ruins and the adult return to Santa Rosa Cave}
\PlanSubsection{Papas' Letter, the Zenithian Sword, and deeper treasure}
\PlanSection{Alcapa revisited: the inn, rumors, and preparations for Reinhart}
\PlanSection{Reinhart Castle basement: entry route, prisoners, and the key}
\PlanSection{Nunnery return and Maria's role at the Southern Tower}
\PlanSection{Southern Tower: floor route, Mirror of Ra, and treasure}
\PlanSection{The False Queen, restored Reinhart, and Henry's departure}
\PlanSection{Reinhart-region cleanup and the next continent}

\PlanChapter{walk-marriage}{Second Generation: Port Selmi and Marriage}
\PlanSection{Port Selmi: the tavern confrontation and Kabochi request}
\PlanSection{Kabochi and the West Kabochi Cave}
\PlanSubsection{Bianca's Ribbon and reunion with the Killer Panther}
\PlanSubsection{Papas' Sword, the village reward, and available recruits}
\PlanSection{Ruraphin: Benett, the Lunar Herb, and learning Return}
\PlanSection{Ruraphin South Cave and the road to Salabona}
\PlanSection{Salabona: Ludman, Flora, and the wedding prerequisites}
\PlanSection{Volcano Cave: route, hazards, treasure, and the Ring of Flames}
\PlanSubsection{Lava Men battle and sustainable healing}
\PlanSection{Mountain Village: Duncan, Bianca, and opening the river route}
\PlanSection{Waterfall Cave: treasure route and the Ring of Water}
\PlanSection{Choosing Bianca or Flora}
\PlanSubsection{Party availability, spell progression, and route differences}
\PlanSubsection{Wedding preparations, Silk Veil, and the ceremony}
\PlanSubsection{Post-wedding gifts, ship access, and what the choice does not change}
\PlanSection{Marriage-region treasure and optional-objective checklist}

\PlanChapter{walk-granvania}{Second Generation: The Road to Granvania}
\PlanSection{The first sea journey and optional coastal stops}
\PlanSection{Medal King Castle, islands, shrines, and travel planning}
\PlanSection{Telepador: desert approach and the Zenithian Helmet test}
\PlanSection{Mountain Cave, west side: ascent, treasure, and resources}
\PlanSection{Chizod / Chizot: the village and the mountain crossing}
\PlanSection{Mountain Cave, east side: descent and hidden treasure}
\PlanSection{Granvania / Granbania: Sancho, Ojiron, family, and castle searches}
\PlanSection{Cave of Trials: puzzles, the Symbol of Royalty, and Kandar}
\PlanSection{Coronation, the twins, and the missing queen}
\PlanSection{Flying Shoes, the Northern Shrine, and the Demon Tower approach}
\PlanSection{Demon Tower: moving floors, treasure, and gatekeeper battles}
\PlanSubsection{Orc and Wyvern gatekeepers}
\PlanSubsection{Jami, the final confrontation, and the generation transition}
\PlanSection{Second-generation completion and carryover checklist}

\PlanChapter{walk-reunion}{Third Generation: Reunion and El Heven}
\PlanSection{The years in stone, rescue, and the return to Granvania}
\PlanSection{The twins, Sancho, Pippin, and rebuilding the active party}
\PlanSection{Zenithian equipment checklist and the son's role}
\PlanSection{Reinhart revisited: Henry, Del, and family-related rewards}
\PlanSection{Northern Church, small shrines, and northeast exploration}
\PlanSection{Water Shrine / Water Cave: the passage to El Heven}
\PlanSection{El Heven: family history, the Magic Key, and the Magic Carpet}
\PlanSection{Newly accessible towns, islands, treasure, and monster recruits}
\PlanSection{Return to Telepador and collection of the Zenithian Helmet}
\PlanSection{Route-order options and prerequisites for restoring Zenithia}

\PlanChapter{walk-zenithia}{Third Generation: Restoring Zenithia}
\PlanSection{Sky Tower: floor route, treasure, and the Magma Staff}
\PlanSection{Opening the mountain barrier and entering Torokko Mines}
\PlanSection{Torokko Mines: trolley switches, loops, Pusan, and the lake}
\PlanSection{The sunken Zenithian Castle and the missing Golden Orb}
\PlanSection{Forest of Illusions: navigating the maze with a child}
\PlanSection{Faerie Village, the Horn of the Fairies, and Faerie Palace}
\PlanSection{The childhood return and the Golden Orb exchange}
\PlanSection{Restoring Zenithia and unlocking castle travel}
\PlanSection{Salabona's sealed jar, Buorn / Bohen, and the Final Key}
\PlanSection{Final Key cleanup and the Royal Cloak side route}
\PlanSection{Tower of Boble: Grappling Hook, dragon statue, and floor route}
\PlanSubsection{Gonz, Gema, and the two Dragon Eyes}
\PlanSubsection{Dragon Orb, Zenithian Armor, and remaining tower treasure}
\PlanSection{The Master Dragon and Zenithian Bells}
\PlanSection{Great Temple: entry, captivity rooms, and party restrictions}
\PlanSubsection{Temple guardians, the high priest, and the wife's rescue}
\PlanSubsection{Ring of Life and preparations for the Demon World}

\PlanChapter{walk-finale}{Third Generation: The Demon World and Ending}
\PlanSection{Water Shrine revisited and the three-ring portal}
\PlanSection{The Demon World overworld: routes, encounters, and recruitment}
\PlanSection{Shahanna: shops, inn, final equipment, and safe preparation}
\PlanSection{Evil Mountain: entrance, summit route, and family reunion}
\PlanSection{Sacred Pitcher, lava removal, floor loops, and navigation puzzles}
\PlanSection{Sage's Stone, Sands of Time, and the final treasure sweep}
\PlanSection{Final guardians and preparation for Mildrath}
\PlanSection{Mildrath's first form: tactics and summoned enemies}
\PlanSection{Mildrath's second form: healing, defenses, and sustained damage}
\PlanSection{Ending tour, credits, clear-state behavior, and postgame access}
\PlanSection{Final-dungeon and pre-ending completion checklist}

\PlanChapter{walk-postgame}{Postgame and the Bonus Dungeon}
\PlanSection{Bonus Dungeon entrance and access prerequisites}
\PlanSection{Floor-by-floor route, branching passages, and puzzle solutions}
\PlanSection{Bonus treasure, the Drum of War, and unique equipment}
\PlanSection{Rare companions and postgame recruitment routes}
\PlanSection{Estark / Esterk: party setup, resistances, healing, and strategy}
\PlanSection{Repeat battles, challenge goals, and remaining completion tasks}

\PlanPart{Characters, Monster Companions, and Abilities}
\PlanChapter{humans}{Human and Temporary Characters}
\PlanSection{Profile fields: availability, role, attributes, spells, and equipment}
\PlanSection{The Hero: childhood, adulthood, and event-learned spells}
\PlanSection{Bianca: childhood carryover and the adult marriage route}
\PlanSection{Flora: second-generation restrictions and later progression}
\PlanSection{The Son: Zenithian equipment, support magic, and combat role}
\PlanSection{The Daughter: offensive magic, support, and equipment}
\PlanSection{Sancho: recruitment, utility, and party role}
\PlanSection{Pippin: recruitment, equipment, and physical combat}
\PlanSection{Temporary and AI-controlled companions}
\PlanSubsection{Papas}
\PlanSubsection{Bella / Vera}
\PlanSubsection{Henry}
\PlanSection{Human spell-learning tables and event-dependent progression}
\PlanSection{Marriage comparison and family availability by generation}
\PlanSection{Supporting characters and relationship glossary}

\PlanChapter{recruitment}{Monster Recruitment and Party Planning}
\PlanSection{Recruitment unlocks, wagon ownership, and eligibility}
\PlanSection{Battle resolution, enemy groups, recruitment chances, and exceptions}
\PlanSection{Guaranteed story companions and random monster recruitment}
\PlanSection{Intelligence, direct commands, and improving obedience}
\PlanSection{Monster Gramps: storage, retrieval, capacity, and releases}
\PlanSection{Repeated species, default names, and roster organization}
\PlanSection{Recruitment routes by generation and region}
\PlanSubsection{Oracleberry, Santa Rosa, and Reinhart}
\PlanSubsection{Port Selmi, Kabochi, Ruraphin, and Salabona}
\PlanSubsection{Telepador, mountain routes, and Granvania}
\PlanSubsection{El Heven, Zenithia routes, and the sea}
\PlanSubsection{Demon World and Bonus Dungeon}
\PlanSection{Party roles: damage, healing, support, resilience, and utility}
\PlanSection{Recommended teams for each story phase and major boss}
\PlanSection{Level caps, training priorities, and late-game replacements}
\PlanSection{Rare recruits, metal companions, and realistic collection planning}

\PlanChapter{companions}{Complete SNES Monster Companion Catalog}
\PlanSection{Companion profile fields}
\PlanSection{Individual monster profiles}
\PlanSubsection{Attackbot}
\PlanSubsection{Big Eye}
\PlanSubsection{Blizzard Hawk}
\PlanSubsection{Bomb Baby}
\PlanSubsection{Bomb Crag}
\PlanSubsection{Brownie}
\PlanSubsection{Centaurus}
\PlanSubsection{Curer}
\PlanSubsection{Dancing Jewel}
\PlanSubsection{Dancing Needle}
\PlanSubsection{Dragon Pup}
\PlanSubsection{Drakee}
\PlanSubsection{Eigerhorn}
\PlanSubsection{Eliminator}
\PlanSubsection{Farewell Crag}
\PlanSubsection{Gigantes}
\PlanSubsection{Golem}
\PlanSubsection{Great Dragon}
\PlanSubsection{Healer}
\PlanSubsection{Hellbattler}
\PlanSubsection{Hork}
\PlanSubsection{Killer Panther: childhood and adult story recruitment}
\PlanSubsection{King Healer}
\PlanSubsection{King Slime}
\PlanSubsection{Kukkle}
\PlanSubsection{Leaonar}
\PlanSubsection{Lionex}
\PlanSubsection{Mad Dragon}
\PlanSubsection{Magician}
\PlanSubsection{Messala}
\PlanSubsection{Metal Babble}
\PlanSubsection{Minidemon}
\PlanSubsection{Neileus}
\PlanSubsection{Orc King}
\PlanSubsection{Puppet Man}
\PlanSubsection{Slime}
\PlanSubsection{Slime Knight}
\PlanSubsection{Soldier Bull}
\PlanSubsection{Springer}
\PlanSubsection{Wyvern}
\PlanSubsection{Yeti}
\PlanSection{Recruitment-location and availability tables}
\PlanSection{Starting abilities, learned abilities, and maximum-level tables}
\PlanSection{Equipment groups, resistance comparisons, and party roles}
\PlanSection{Companion names and aliases}

\PlanChapter{abilities}{Spells and Special Abilities}
\PlanSection{Alphabetical index: MP cost, targets, effects, and users}
\PlanSection{Healing, revival, sacrifice, and status recovery}
\PlanSection{Offensive magic by family and target pattern}
\PlanSection{Instant death, sleep, confusion, silence, and other control effects}
\PlanSection{Attack and defense changes, barriers, reflection, and dispelling}
\PlanSection{Breath attacks, dances, physical skills, and monster-only abilities}
\PlanSection{Travel and field utility: Return, Outside, Stepguard, and time changes}
\PlanSection{Event-learned spells: Return and Chance}
\PlanSection{Learning-level tables for humans and monster companions}
\PlanSection{Item-cast equivalents, AI use, resistance, and reliability}

\PlanPart{Items, Equipment, Enemies, and Bosses}
\PlanChapter{items}{Items, Keys, and Story Objects}
\PlanSection{Alphabetical item index and acquisition cross-reference}
\PlanSection{Healing, status cures, revival, and MP-restoration consumables}
\PlanSection{Seeds, acorns, permanent growth, and allocation planning}
\PlanSection{Battle items, reusable effects, limited uses, and breakage}
\PlanSubsection{Sage's Stone, Drum of War, and Sands of Time}
\PlanSubsection{Rings, stones, waters, bells, and situational tools}
\PlanSection{Keys and access progression}
\PlanSubsection{Thieves' Key knowledge, Reinhart Key, Magic Key, and Final Key}
\PlanSection{Childhood and faerie-realm objects}
\PlanSubsection{Torch, Golden Orb, Flute of Spring, and Bianca's Ribbon}
\PlanSection{Wedding, royal, and rescue objects}
\PlanSubsection{Papas' Letter, Mirror of Ra, wedding rings, and Silk Veil}
\PlanSubsection{Symbol of Royalty, Flying Shoes, and Ring of Life}
\PlanSection{Zenithia and Demon World objects}
\PlanSubsection{Magma Staff, faerie horn, orbs, Grappling Hook, and Dragon Eyes}
\PlanSubsection{Dragon Orb, Zenithian Bells, and Sacred Pitcher}
\PlanSection{Maps, wagon, vehicles, and non-inventory unlocks}
\PlanSection{Where to obtain, retain, use, or revisit each key object}

\PlanChapter{equipment}{Weapons, Armor, Shields, Helmets, and Accessories}
\PlanSection{Catalog fields: power, prices, locations, users, and special effects}
\PlanSection{Complete weapon catalog}
\PlanSubsection{Single-target weapons and character restrictions}
\PlanSubsection{Whips, boomerangs, repeated attacks, and unusual damage effects}
\PlanSubsection{Weapon-cast spells and reusable support tools}
\PlanSection{Complete body-armor catalog}
\PlanSection{Complete shield catalog}
\PlanSection{Complete helmet catalog}
\PlanSection{Complete accessory catalog}
\PlanSection{Human and monster equipment-compatibility matrices}
\PlanSection{Cursed equipment, drawbacks, removal, and sale restrictions}
\PlanSection{Elemental protection, status protection, and special properties}
\PlanSection{Zenithian equipment: acquisition, eligibility, and checklist}
\PlanSection{Casino prizes, Medal King rewards, unique finds, and monster drops}
\PlanSection{Recommended upgrades by generation, role, and budget}
\PlanSection{Endgame and postgame loadouts for each party member}

\PlanChapter{bestiary}{Enemy Bestiary, Encounters, and Treasure Drops}
\PlanSection{Bestiary fields: names, HP, rewards, drops, behavior, and locations}
\PlanSection{Alphabetical enemy catalog}
\PlanSubsection{Enemies A--F}
\PlanSubsection{Enemies G--L}
\PlanSubsection{Enemies M--R}
\PlanSubsection{Enemies S--Z}
\PlanSection{Overworld encounter tables by region and generation}
\PlanSection{Dungeon encounters by floor and route}
\PlanSection{Sea encounters and vehicle-dependent access}
\PlanSection{Metal enemies: locations, escape behavior, and training tactics}
\PlanSection{Mimics, trapped chests, summoned enemies, and fixed encounters}
\PlanSection{Drop tables and item-to-monster reverse index}
\PlanSection{Recruitable flags and differences between enemy and companion data}
\PlanSection{Boss entries and variant encounters}

\PlanChapter{bosses}{Boss and Scripted-Battle Strategy Compendium}
\PlanSection{Strategy format: access, party, equipment, patterns, and recovery}
\PlanSection{First-generation battles}
\PlanSubsection{Lenoire ghost boss}
\PlanSubsection{Zaile and the Ice Mansion boss}
\PlanSubsection{Gema and the childhood scripted confrontation}
\PlanSection{Second-generation battles}
\PlanSubsection{False Queen}
\PlanSubsection{Killer Panther reunion: resolution without defeating it}
\PlanSubsection{Lava Men}
\PlanSubsection{Kandar and the Cave of Trials confrontation}
\PlanSubsection{Demon Tower gatekeepers and Jami}
\PlanSection{Third-generation battles}
\PlanSubsection{Buorn / Bohen}
\PlanSubsection{Gonz and Gema at the Tower of Boble}
\PlanSubsection{Great Temple guardians and high priest}
\PlanSection{Final-dungeon encounters and both Mildrath forms}
\PlanSection{Estark / Esterk and repeat postgame challenges}
\PlanSection{Low-resource recovery plans and common causes of failed attempts}

\PlanPart{Optional Content, Economy, and Completion}
\PlanChapter{medals}{Small Medals and Medal King Rewards}
\PlanSection{Medal collection, exchange rules, rewards, and planning}
\PlanSection{Chronological location checklist}
\PlanSubsection{Early adulthood and the Reinhart region}
\PlanSubsection{Port Selmi, Kabochi, Ruraphin, and Salabona}
\PlanSubsection{Mountain routes, Granvania, Cave of Trials, and Demon Tower}
\PlanSubsection{Sea routes, shrines, El Heven, and key-gated revisits}
\PlanSubsection{Zenithia routes, Tower of Boble, and Great Temple}
\PlanSubsection{Demon World, Evil Mountain, and Bonus Dungeon}
\PlanSection{Searchable objects, fixed rewards, and enemy-drop medals}
\PlanSection{Location-to-map index and the final collection audit}
\PlanSection{Version-specific Mimic encounters}

\PlanChapter{economy}{Casino, Shops, and Resource Planning}
\PlanSection{Oracleberry casino access, coins, stakes, and prize exchange}
\PlanSection{Slime races}
\PlanSection{Slot machines}
\PlanSection{Monster arena betting}
\PlanSection{Prize catalog, acquisition priorities, and equipment tradeoffs}
\PlanSection{Ordinary-play budgeting versus optional emulator-assisted methods}
\PlanSection{Town-by-town weapon, armor, and item shop inventories}
\PlanSection{Day/night availability, marriage-route support, and special vendors}
\PlanSection{Gold, experience, healing costs, and efficient purchase routes}
\PlanSection{Banks, retained valuables, sale decisions, and late-game funding}

\PlanChapter{optional}{Secrets, Side Quests, and Optional Objectives}
\PlanSection{Reinhart revisits, family rewards, and travel privileges}
\PlanSection{Ruraphin follow-up, Benett's experiments, and Chance}
\PlanSection{Medal King Island, optional sea stops, and newly opened doors}
\PlanSection{Buorn / Bohen: timing, sealed-jar events, and Final Key access}
\PlanSection{Emblem Cave / Cave of Seals and the Royal Cloak}
\PlanSection{Hidden mountain-path treasure and the Staff of Punishment}
\PlanSection{Optional recruits, unique equipment, and reusable support items}
\PlanSection{Inventory-order Easter egg: requirements and effects}
\PlanSection{Version-specific Mimic encounters}
\PlanSection{Emulator save-state casino methods, clearly separated from normal play}
\PlanSection{Postgame secrets and version-specific misconceptions}

\PlanChapter{completion}{Completion Checklists and Efficient Routes}
\PlanSection{Before leaving childhood: treasures, companions, and irreversible events}
\PlanSection{Before marriage and before the second-generation transition}
\PlanSection{Third-generation unlock and backtracking checklist}
\PlanSection{Every town and dungeon treasure checklist}
\PlanSection{All Small Medals and reward selections}
\PlanSection{All recruitable monster types and story companions}
\PlanSection{All unique equipment, key objects, and reusable battle tools}
\PlanSection{Spell progression, event spells, and companion level-cap goals}
\PlanSection{Before Mildrath and before entering the Bonus Dungeon}
\PlanSection{Postgame completion, repeat challenges, and optional personal goals}
\PlanSection{Suggested save checkpoints and recovery routes}

\PlanPart{World, Town, and Dungeon Atlas}
\PlanChapter{world-atlas}{World Map, Towns, Services, and Travel Routes}
\PlanSection{Overworld map, region names, and location index}
\label{map-overworld}
\PlanSection{Travel-unlock and access-dependency diagrams}
\PlanSection{Ports, landings, Return points, carpet limits, and dragon access}
\PlanSection{Town and settlement profiles}
\PlanSubsection{Vista Port, Santa Rosa, and Alcapa / Alkapa}
\PlanSubsection{Reinhart, Nameless Nunnery, and Oracleberry}
\PlanSubsection{Port Selmi, Kabochi, and Ruraphin}
\PlanSubsection{Salabona and Mountain Village}
\PlanSubsection{Telepador, desert stops, and Medal King Castle}
\PlanSubsection{Chizod / Chizot, mountain stops, and Granvania}
\PlanSubsection{Northern Church, small shrines, and El Heven}
\PlanSubsection{Faerie Village, Faerie Palace, and Zenithian Castle}
\PlanSubsection{Shahanna and the Demon World}
\PlanSection{Service matrix: shops, inns, churches, banks, and monster storage}
\PlanSection{Generation changes, day/night differences, and key-gated searches}

\PlanChapter{dungeon-atlas}{Dungeon Atlas, Puzzles, and Treasure Routes}
\PlanSection{Atlas conventions: floors, entrances, exits, switches, and treasure}
\PlanSection{First-generation maps}
\PlanSubsection{Santa Rosa Cave, including the adult revisit}
\label{map-santa-rosa}
\PlanSubsection{Lenoire / Linoire Castle}
\label{map-linoire}
\PlanSubsection{Dwarves Cave}
\label{map-dwarves}
\PlanSubsection{Ice Mansion}
\label{map-ice}
\PlanSubsection{Thieves Cave / Reinhart Cave}
\label{map-thieves}
\PlanSection{Reinhart, regional side quests, and marriage maps}
\PlanSubsection{Reinhart Castle basement}
\label{map-reinhart}
\PlanSubsection{Southern Tower}
\label{map-southern}
\PlanSubsection{Ruraphin South Cave}
\label{map-ruraphin}
\PlanSubsection{West Kabochi Cave}
\label{map-kabochi}
\PlanSubsection{Volcano Cave}
\label{map-volcano}
\PlanSubsection{Waterfall Cave}
\label{map-waterfall}
\PlanSection{Mountain crossing, royal trials, and rescue maps}
\PlanSubsection{Mountain Cave, west side}
\label{map-mountain-west}
\PlanSubsection{Mountain Cave, east side}
\label{map-mountain-east}
\PlanSubsection{Cave of Trials}
\label{map-trials}
\PlanSubsection{Demon Tower}
\label{map-demon-tower}
\PlanSection{Third-generation travel, Zenithia, and temple maps}
\PlanSubsection{Water Shrine: El Heven route and ring portal}
\label{map-water-shrine}
\PlanSubsection{Emblem Cave / Cave of Seals}
\label{map-emblem}
\PlanSubsection{Sky Tower}
\label{map-sky}
\PlanSubsection{Torokko Mines / Trolley Cave}
\label{map-torokko}
\PlanSubsection{Forest of Illusions}
\label{map-forest}
\PlanSubsection{Tower of Boble}
\label{map-boble}
\PlanSubsection{Great Temple}
\label{map-great-temple}
\PlanSection{Final and postgame dungeon maps}
\PlanSubsection{Evil Mountain}
\label{map-evil-mountain}
\PlanSubsection{Bonus Dungeon}
\label{map-bonus}
\PlanSection{Treasure, Small Medal, key-object, and encounter overlays}
\PlanSection{Text routes and additional locations}

\PlanPart{Appendices}
\PlanChapter{names}{Name Cross-Reference and Glossary}
\PlanSection{People, family relationships, and temporary companions}
\PlanSection{Towns, dungeons, regions, and vehicles}
\PlanSection{Monster species and default companion names}
\PlanSection{Spell names and aliases}
\PlanSection{Items, equipment, key objects, and naming collisions}
\PlanSection{Alternate names for lookup}

\PlanChapter{quick-reference}{Quick Reference Tables}
\PlanSection{Story milestones and next-objective lookup}
\PlanSection{Keys, travel unlocks, and access requirements}
\PlanSection{Human and monster spell-learning summaries}
\PlanSection{Monster level caps, recruitment regions, and party roles}
\PlanSection{Equipment compatibility, upgrades, and resistance priorities}
\PlanSection{Boss preparation, status counters, and reusable-item lookup}
\PlanSection{Small Medal rewards, casino prizes, and high-value purchases}
\PlanSection{Shops, safe recovery points, and final completion checklists}

\tableofcontents

\clearpage
\begingroup
\renewcommand{\addcontentsline}[3]{}
\renewcommand{\theHchapter}{written.\arabic{chapter}}
\setcounter{chapter}{0}
\chapter{Game Version, Names, and Guide Conventions}
\label{written-conventions}

\emph{Dragon Quest V} begins aboard a ship, with a young boy travelling
alongside his father, Papas. The first visits to Santa Rosa and Alcapa
establish a comfortable rhythm of talking to townspeople, searching houses,
buying equipment, and venturing into the countryside. As the journey grows,
the party gains new companions and new ways to travel. Familiar places
remain part of the adventure, and return visits give them fresh significance.

The game rewards attention to small details. A drawer can hold a useful
healing item. A conversation can point toward the next destination. A new
key can make several earlier stops worth revisiting. Keeping track of these
connections makes exploration satisfying and helps the party arrive ready
for the next challenge.

\section{Game Version}

\emph{Dragon Quest V: Tenkuu no Hanayome} is Enix's 1992 role-playing game
for the Super Famicom / SNES.

The subtitle is commonly rendered as \emph{The Bride of the Sky} or
\emph{Bride of Heaven}. The names \emph{Dragon Quest V}, \emph{Dragon
Quest 5}, and \emph{DQ5} identify the same game.

Exploration, party development, and monster recruitment are central to the
adventure. Several features shape the journey:
\begin{itemize}
  \item The normal active battle party has three members. Choosing their
    roles and equipment is an important part of preparing for a dungeon.
  \item The wagon carries additional companions. Access to those companions
    depends on the location and the party's current circumstances.
  \item Each recruitable monster species has its own growth, abilities,
    equipment options, and maximum level.
  \item Keys and vehicles expand the accessible world throughout the game.
    A new travel ability can make an earlier destination worth another visit.
\end{itemize}

\section{Names}

\subsection*{One place, several names}

Alcapa, Alkapa, and Arukapa identify the same town. Lenoire and Linoire
are alternate spellings of the same castle. Common place names include
the following:

\begin{center}
\begin{tabular}{@{}p{0.28\textwidth}p{0.64\textwidth}@{}}
\textbf{Primary name} & \textbf{Alternate names} \\
\hline
Santa Rosa & Santa Rose, Santorez \\
Alcapa & Alkapa, Arukapa \\
Lenoire Castle & Linoire Castle, Lenule, Linule \\
Oracleberry & Oracle Valley \\
Salabona & Sarabona \\
Granvania & Granbania, Glanvania \\
El Heven & Elheaven \\
Torokko Mines & Trolley Cave \\
\end{tabular}
\end{center}

Surroundings provide another useful way to identify a place. Its position
relative to a port, mountain pass, river, or nearby town helps confirm the
match. Dungeon landmarks, floor connections, and the current objective
provide similar clues when a place has an unfamiliar name.

\subsection*{People and companion names}

The \emph{Hero} is the player-named protagonist who travels with Papas at
the beginning of the adventure.

Monster companions have both a species and an individual name. A Slime
Knight's species determines its growth, spells, and equipment options.
Individual companions can have their own nicknames, and the roster may
contain more than one member of a species.

A character's appearance, location, and abilities help establish its
identity. These details are useful when several enemy or companion names
resemble one another.

\subsection*{Spells and equipment}

\emph{Blaze}, \emph{Firebal}, and \emph{Bang} belong to separate spell
families. MP cost, target pattern, and effect help identify a spell and
its role in battle.

Weapons and armor have individual names, power values, eligible users,
and special properties. Certain rings and staves also have item-use
effects alongside their equipment properties.

\chapter{Controls, Exploration, and Progression}
\label{written-exploration}

A visit to town gives the party a chance to recover, hear local news, and
prepare for the next trip. Out in the countryside, the same small decisions
continue to matter. Healing before a difficult stretch, keeping useful
items with the right companion, and remembering a convenient resting place
can make a long journey considerably more comfortable.

\section{Controller inputs, field commands, and menu navigation}

\subsection*{Moving and choosing an action}

Use the directional pad to walk and to move the cursor through menus.
Approach a person or object and face it before choosing an interaction.
A staircase, entrance, or passage generally connects to the next area when
the party walks onto the appropriate square.

Press \textbf{A} to open the field command menu. Choose a command with the
directional pad and confirm it with A. In menus, \textbf{B} cancels the
current selection or closes the current window. Some actions lead through
several choices, such as selecting a companion, an item, and its recipient.
Read the current prompt before confirming the next choice.

\subsection*{The field commands}

\begin{description}
  \item[Talk:] Speak to the person the leading character is facing. Shop
    counters and service desks are approached in the same way.
  \item[Spell:] Choose a companion and one of that companion's field spells.
    Healing, status recovery, and travel spells appear as they become
    available. Select the recipient when the spell requests a target.
  \item[Status:] Inspect a companion's level, attributes, HP, MP, equipment,
    and other personal information. Use this screen when comparing an
    upgrade or checking the party's condition before leaving town.
  \item[Item:] Choose an owner and an object from that owner's inventory.
    The available actions include using, transferring, and discarding
    ordinary items. An item's available actions depend on the object.
  \item[Equip:] Select a companion and assign suitable weapons, armor,
    shields, and headgear. Check the resulting attack and defense values
    after making a change.
  \item[Tactics:] Manage party-related commands and battle behavior.
    Party arrangement determines the companions travelling at the front.
  \item[Door:] Open the door directly ahead when the party has the required
    access. Approach the door square and face it before using the command.
  \item[Search:] Examine the current position or the object directly ahead.
    This command is useful for chests, drawers, pots, and specific places
    on the ground.
\end{description}

A character-selection window identifies the companion whose spell,
equipment, or inventory is being used. Confirm the owner carefully when
several companions carry similar objects. Giving a Medical Herb to a
front-line companion places that herb among the items that companion can
use during a battle.

\subsection*{Working through menus}

Shops, inns, churches, and storage counters lead through a short sequence
of questions. Select the service, choose the companion or object involved,
and review the price or destination before accepting. A purchase may also
ask which companion should receive the item.

When the party has accumulated several spells or items, move through the
list to inspect the available choices. Returning to the previous window
lets you choose another owner or action. This is useful when distributing
new equipment or preparing several companions for a dungeon.

The opening data menu handles starting an adventure and managing its save
slots. Copying creates another copy of the selected adventure record.
Deleting removes the chosen record. Check the displayed slot and character
name carefully before confirming either operation.

\section{Data and status screens: attributes, experience, and equipment}

\subsection*{Checking a companion}

The Status command gives a compact view of a companion's development and
current condition. Begin with the name and level, then check HP and MP.
Their current values show the resources available for the next trip, and
their maximum values show the companion's capacity when fully recovered.

A character with ample maximum MP can still be nearly exhausted after a
long dungeon. Check the current value before relying on another healing
spell or planning a magical escape. The same habit helps identify a
companion who needs treatment before the next battle.

Experience records progress toward the next level. Level gains can improve
attributes and add spells or abilities. Check the Spell menu after a useful
level gain to see which new field abilities are available. A priest's
experience advice gives the amount still needed for the next level.

\subsection*{Attributes and equipment}

Strength, agility, base defense, intelligence, and luck describe the
companion's current development. Strength contributes to physical attack
power. Agility influences the order of actions in battle. Base defense
contributes to protection from physical attacks. Intelligence affects a
monster companion's command obedience.

Attack and defense show the current results of character development and
equipment. When considering a purchase, inspect the intended user's
existing weapon or armor first. After equipping the replacement, check
the new total and decide where the old piece will be most useful.

Special equipment properties deserve a separate look. Protection from
certain attacks or status effects can be valuable in the next dungeon.
Some weapons have unusual attack patterns, and some objects provide a
spell-like effect when used as an item. Choose equipment with the area's
enemies and the companion's role in mind.

\subsection*{Growth and the active party}

Monster species have individual level caps. A companion approaching its
cap may already have most of its available abilities. Consider its healing,
support, equipment, and durability when deciding how long to keep it in
the active party.

After a recruitment or party change, inspect the new companion's status,
spells, and equipment. Supply suitable gear and place useful consumables
in its inventory. This makes the new member ready to contribute on the
next journey.

The opening data menu also provides message-speed and sound settings.
Choose a comfortable text speed so that shop prices, service prompts,
and conversations remain easy to follow.

\section{Dialogue, searching, drawers, pots, chests, and hidden objects}

\subsection*{Talking through a town}

Speak to residents as you explore. Their conversations introduce nearby
places, local problems, unusual visitors, and changes in the town. A
shopkeeper may offer a useful destination, while another resident explains
why a building or route matters.

Conversations can change after an event or at another time of day. Return
to the people involved in the current objective after completing its
immediate task. The next instruction, reward, or permission to proceed
often comes through that follow-up conversation.

Face the person before using Talk. Across a counter, face the clerk from
the customer side. When a person has moved to another room or a building
has closed, check the time of day and look around the surrounding area.

\subsection*{Searching houses and shops}

Drawers, pots, barrels, and chests can hold useful objects. Approach and
face the container, then use Search. Check accessible rooms above or
behind shops as well as the obvious ground-floor rooms. A small find can
save a purchase or supply an item for the next trip.

The opening ship offers an immediate example. Search the drawer in the
cabin for a Medical Herb. In Santa Rosa, household searches provide early
supplies and equipment. These discoveries make it worth examining a new
room before leaving it.

Keep some inventory space available while exploring. Transfer spare
objects between companions or deposit them before a long treasure run.
A party carrying a full assortment of old equipment can run short of
space just as a useful new object is found.

\subsection*{Ground searches and unusual containers}

Some objects lie on a particular square of ground. Stand on the indicated
square and use Search. Details such as a flower patch, grave, well,
bucket, or conspicuous corner can help identify the correct position.
NPC dialogue may point toward the search area.

During a return visit, check rooms that have become accessible through
a new key or event. A newly opened cellar or side room may contain several
objects close together. Finish the searches in that room before returning
to the main route.

Some containers begin a fixed encounter. Keep the party reasonably healed
while opening unfamiliar containers in a dungeon, and have a recovery
item ready in an active companion's inventory.

\section{Day and night: access, shops, events, and changing time}

\subsection*{A town at two different times}

Day and night affect the daily life of towns. Residents move, businesses
open or close, and certain people offer their services at a particular
time. The change in lighting also makes it easy to recognize which part
of the day the party has reached.

Oracleberry provides a useful early example. Its special merchant conducts
business at night, and the fortune teller also has a nighttime routine.
Visit the town during both periods to see the available services and
conversations.

When a person is absent or a counter is closed, check nearby rooms and
consider another time of day. A resident may explain the schedule. Some
story events have their own timing and require a conversation or action
before the party can continue.

\subsection*{Letting time pass}

Travel around the countryside until the lighting changes, then enter
the town at the desired time. Choose a manageable area near the entrance
when waiting for night. Encounters during the walk give the party a chance
to earn experience and gold, so keep healing supplies available.

An ordinary stay at an inn restores the party and brings the day around
to morning. A story-directed stay may include an event. Once the party
is free to move again, check the surroundings and continue the current
objective.

\subsection*{Changing the time directly}

The Day-Night spell costs 3 MP and switches between day and night when
used in an appropriate field location. Select a companion who knows it
through the Spell command. This becomes convenient when returning to a
shop or searching for a person with a particular schedule.

The Lamp of Darkness changes day into night. Use it through the owner's
Item menu when that effect is needed. Enter or revisit the town afterward
to reach the nighttime activity.

\section{Doors, keys, moving obstacles, floor hazards, and puzzles}

\subsection*{Opening a route}

Approach a closed door and face it, then use Door. Locked doors require
the relevant key or learned unlocking ability. Keep important access
items available to the travelling party.

The Thieves' Key knowledge opens the ordinary locks associated with that
ability. The Reinhart Key is used for the castle's relevant locks. The
Magic Key opens its red-framed doors, and the Final Key provides access
to prison cells and other remaining key-locked doors.

A new key can make several older destinations worth revisiting. Complete
the newly accessible searches in one town before moving on to the next.
An entrance controlled by a story event or a mechanism requires that
particular action as well.

\subsection*{Rocks, switches, and connected passages}

Some movable obstacles can be pushed by walking into them from a suitable
side. Check the clear space around the object before moving it. A passage
can depend on where an obstacle ends up, so pay attention to nearby walls,
doors, and floor connections.

Inspect visible switches, levers, and unusual structures. Look for the
gate, track, staircase, or passage affected by the mechanism. After an
action, examine the changed area before moving farther away.

In dungeons with holes, sliding floors, or one-way connections, keep track
of the stairs and entrances that reconnect the rooms. A fall can lead to
a different section of the same dungeon. Locate the return path before
spending too many healing resources exploring its branches.

\subsection*{Crossing damaging terrain}

Hazardous floors can reduce HP while the party moves. Check the party's
condition after crossing them, especially when a battle may begin nearby.
Poison can also drain HP during travel and deserves prompt treatment.

Stepguard costs 2 MP and protects the party from damaging floor tiles.
Cast it before crossing an appropriate hazard. Check the party's condition
as the route moves into another area, and renew the protection when needed.

If the caster is short of MP, choose a route with fewer hazardous steps
and heal after the crossing. Keeping the party healthy before opening a
door or entering a new room gives more room to handle an unexpected battle.

\section{Inventory ownership, item transfer, capacity, and equipment use}

\subsection*{The twelve-item inventory}

Each companion has twelve inventory slots. Equipment occupies those
slots along with consumables, spare gear, and other carried objects.
Separate copies of a consumable occupy separate slots.

Leave room for discoveries before entering a treasure-rich area. Deposit
unused possessions, sell ordinary spare gear when it is safe to do so,
and share the carrying load among available companions. Keep the active
party's recovery supplies convenient to reach.

Item ownership matters in battle. A companion uses objects from its own
inventory during its turn. Give healing or status-recovery items to active
members who can use them when needed. A companion with a useful reusable
item should also be available for the encounter where it will be used.

\subsection*{Transferring and equipping}

Open Item, select the current owner, choose the object, and select the
transfer action. Choose the recipient and confirm the transfer. Check
the recipient's remaining space when moving several objects at once.

Use Equip to assign a suitable weapon, armor piece, shield, or headgear.
After a change, inspect the companion's attack and defense values. The
old equipment remains a possession until it is transferred, sold, stored,
or otherwise removed.

An older weapon or armor piece may suit another companion. Check the
recipient's equipment options before selling it. When moving equipped
gear, review both companions' equipment afterward so that each has the
intended loadout.

\subsection*{Using objects and conserving valuables}

Select Use from the object's available item actions. The game may ask
for a recipient or apply the effect to the field. Medical Herbs provide
healing, Antidotes treat poison, and Full Moon Herbs treat paralysis.
Their availability can help preserve a healer's MP during a long trip.

Certain weapons and other objects produce a spell-like effect when used.
Their useful action may be available alongside their equipment function.
Inspect the object's effect and give it to a companion who can make good
use of it.

Keep important keys, event objects, and unique possessions safe. Store
objects that are temporarily unnecessary when the depository accepts
them. Resolve cursed equipment through the appropriate church service
before trying to change an affected loadout.

\section{Inns, churches, saving, revival, and status recovery}

\subsection*{Resting before the next journey}

An inn restores the living party's HP and MP. Speak to the innkeeper,
review the cost, and accept the stay. Free resting opportunities are
particularly useful during early training or while gathering enough gold
for an equipment upgrade.

Check the party's condition after resting. Treat persistent ailments and
revive fallen companions before setting out. A fully supplied healer and
healthy front line make the next journey much more comfortable.

\subsection*{Church services}

Speak to the priest and select the service needed:
\begin{description}
  \item[Record Adventure:] Save the current adventure. Confirm the chosen
    record and wait for the priest to finish the saving dialogue.
  \item[Experience advice:] Learn how much experience a companion needs
    to reach the next level.
  \item[Revival:] Restore a fallen companion to life. Select the companion
    and review the required payment.
  \item[Poison treatment:] Remove poison from the selected companion.
  \item[Curse removal:] Resolve a curse affecting a companion or its
    equipment through the priest's service.
\end{description}

Save after an important acquisition, a difficult dungeon, or a substantial
party change. Before a longer expedition, visit the inn and church so
that the saved record reflects the prepared party.

\subsection*{Recovery away from town}

Healing spells and consumables let the party recover between battles.
Antidote costs 2 MP and treats poison. Numboff costs 2 MP and treats
paralysis. Carrying Antidotes and Full Moon Herbs gives additional options
when the appropriate caster is unavailable or short of MP.

Vivify costs 10 MP and attempts to revive a fallen ally. A failed attempt
requires another cast, so allow enough MP when relying on it. Revive costs
20 MP and restores a fallen ally with full HP. A Leaf of the World Tree
also provides revival and is a valuable emergency possession.

After a revival, check HP, MP, and remaining supplies. A recovery spell
can make it possible to finish a short route or reach the exit. A return
to town provides a comfortable chance to rest, replenish items, and save.

\subsection*{Recovering from defeat}

When the party is defeated, half of the gold being carried is lost.
Banked gold remains protected. Check the party's condition after returning
to the recovery point, then arrange revival, treatment, and rest.

Revisit the equipment and supply plan before trying the area again.
A protective item, a different active companion, or an extra recovery
object may give the party the help it needs on the next attempt.

\section{Banks, gold storage, shopping, and selling}

\subsection*{Depositing gold and possessions}

The depository offers storage for gold and items. Oracleberry has one
near Monster Gramps. Gold deposits are made in increments of 1,000 G.
Keep enough spending money for the next inn stay, treatment, and planned
purchases, then deposit the amount that can be set aside.

Stored gold is protected from the loss caused by defeat. Withdraw it when
an equipment purchase or recovery expense is due. Keeping a reserve in
the bank makes rebuilding a defeated party easier.

The item-storage service is useful for spare equipment and possessions
that are being kept for later. Select the owner and object when depositing,
then use the withdrawal service when the object is needed again. Leave
space with the intended recipient before collecting several stored items.

When a monster is left with Monster Gramps, its carried equipment and
items are transferred to the depository. Retrieve the needed possessions
and equip the companion again when bringing it back into the travelling
party.

\subsection*{Making a purchase}

Speak to the shopkeeper, choose the item, and review its price and the
intended recipient. Consider which active companion gains the most from
the upgrade. A dependable weapon for a frequent attacker or protective
gear for a vulnerable member can have an immediate effect on the next trip.

After buying equipment, assign it through Equip and check the resulting
status values. Transfer useful older gear to another eligible companion.
Keep enough cash for recovery and essential consumables after the purchase.

Oracleberry's nighttime merchant sells the wagon for 300 G. Having the
required cash ready makes that visit convenient. Purchases tied to travel
or party access are worth accounting for alongside ordinary equipment.

\subsection*{Selling and keeping useful objects}

The selling service quotes a value before the transaction is completed.
Review the object and its owner carefully. Ordinary gear that has been
replaced throughout the party can provide money and free inventory space.

Keep unique objects, useful reusable items, and equipment intended for a
companion who will be trained later. The depository provides space for
possessions that are worth retaining. Store them together with a clear
idea of which companion or journey they will support.

\section{Travel progression: walking, wagon, ship, carpet, castle, and dragon}

\subsection*{Walking and the wagon}

On foot, follow paths through the countryside and look for the entrances
to towns, caves, towers, and other destinations. Mountains and other
obstacles define the available routes. Encounters along the way consume
healing supplies and MP, so consider the journey to the entrance when
preparing for a dungeon.

The wagon follows the party during appropriate field travel and provides
room for additional companions. Arrange the active members and distribute
their equipment before entering an area with restricted wagon access.
Keep the supplies needed inside that area with the companions who will
be available there.

\subsection*{Sailing}

The ship opens sea routes between continents and islands. Board at its
docking point and sail toward the desired coast. Choose an accessible
landing place near the destination, then continue on foot.

Sea travel has its own encounters. Maintain enough healing for the voyage
and the exploration after landing. A nearby town or inn can provide a
convenient starting point for a longer island trip.

Remember the landing point before entering a dungeon or making a detour.
The coastline and nearby landmarks help locate the ship again when the
party is ready to continue.

\subsection*{Using the Magic Carpet}

Use the Magic Carpet from its owner's inventory on suitable open ground.
It carries the party across plains and water, making several routes more
direct. Forests and mountains constrain its movement.

Choose a clear approach to the destination and a suitable place to land.
After dismounting, walk into the town or dungeon. While travelling by
carpet, X opens the world map and helps identify nearby land and routes.

\subsection*{The flying castle}

Once Zenithia is available for travel, its movement opens another way to
cross the world. Use the directional pad while controlling the castle.
Press X to view the world map and check the destination before choosing
the approach.

The castle's interior also provides a place to regroup during that part
of the adventure. Keep the party ready for the destination reached after
leaving it, especially when continuing directly into a dungeon.

\subsection*{Dragon travel}

The Zenithian Bells summon the flying ride when they are available.
Use the item outdoors and choose the desired route through the world.
This form of flight provides access across terrain such as forests and
mountains.

Choose a suitable landing point near the destination, then dismount and
continue on foot. Check the active party and healing supplies before
entering the next area. A fast flight can place the party close to a
demanding dungeon with little travel in between.

\section{Return destinations, Outside, maps, and safe backtracking}

\subsection*{Return and Outside}

Return costs 8 MP and transports the party to an available destination
in its list. Towns and other eligible stops become available through the
journey. Select the desired destination from the Spell command while
standing in a suitable outdoor location.

Outside costs 8 MP and takes the party out of a dungeon that allows
magical escape. It is useful after collecting treasure or when healing
resources are running low. Once outside, the party can choose the next
overland route or use Return.

If the same companion will cast both spells, keep at least 16 MP available
for that trip home. Allow additional MP for any healing still needed before
reaching the exit. The remaining value on the Status screen is useful
when deciding whether to explore another branch.

\subsection*{Other travel aids}

A Wing of Wyvern provides an item-based journey back to the last town
visited. Carry one with an available companion and use it from a suitable
outdoor position when that destination is useful.

Repel costs 4 MP and helps avoid random battles with weak enemies. Holy
Water also helps ward off weak encounters. These aids can make a familiar
journey more convenient as the party grows stronger.

Use the Wayfarer's Map to examine the overworld when the item is available.
Compare the party's position with the coast, rivers, mountains, and nearby
settlements. The map is particularly useful when a new vehicle opens a
route across water or around a large obstacle.

\subsection*{Returning for earlier discoveries}

A new key or travel ability can open several useful return visits. Choose
one town or region, collect its newly accessible treasure, and then move
to the next nearby stop. Visit shops that have become useful to the
current party and retrieve any stored equipment needed for the trip.

Keep a short mental route through the area: the landing or Return point,
the rooms to search, the services to use, and the way back to the vehicle
or next destination. This makes a return visit comfortable even when the
town was last explored much earlier in the adventure.

Finish a longer outing by treating ailments, restoring HP and MP,
reorganizing useful possessions, and recording the adventure at a church.
The party is then ready to begin its next journey from a convenient stop.

\chapter{Combat, Tactics, and Party Management}
\label{written-combat}

Three active companions give every action a noticeable place in a battle.
A timely heal can keep an attacker on its feet, a group spell can clear a
dangerous cluster, and a defensive action can buy the party the round it
needs to recover. The available companions and their equipment give the
player several ways to handle the same encounter.

Watch the enemy groups, the party's condition, and the resources needed
for the rest of the journey. A comfortable victory leaves enough HP, MP,
and useful items to continue exploring. Against a stronger opponent, the
same habits help the party settle into a sustainable sequence of attacks,
support, and recovery.

\section{Battle commands, target groups, turn order, and escape}

\subsection*{The opening battle menu}

The first battle menu offers four main actions:
\begin{description}
  \item[Fight:] Begin choosing the party's actions for the round. The
    current tactics determine which companions receive direct orders.
  \item[Change:] Adjust the active members when reserve companions and
    the encounter's conditions allow it.
  \item[Tactics:] Choose the battle behavior of computer-controlled
    companions or use the direct-command setting.
  \item[Run:] Attempt to leave an encounter that permits escape.
\end{description}

For a companion receiving direct orders, the available commands include
Attack, Spell, Parry, and Item. Attack uses the equipped weapon. Spell
opens the companion's battle abilities. Parry is the defensive action,
reducing incoming damage during that turn. Item opens the objects carried
by that companion.

Select the action and its target, then continue through the other directly
controlled members. Check the party's HP while choosing the round. A
companion who is already badly injured may need healing or a defensive
action before taking another heavy hit.

\subsection*{Reading the enemy groups}

The enemy display identifies the species present and the number remaining
in each group. The target menu uses those entries when selecting an
offensive action. A single-target attack affects one creature in the
chosen group, a group attack affects that group, and an all-enemy action
reaches across the opposing formation.

Choose targets according to their immediate threat. An enemy that heals,
calls reinforcements, drains MP, or inflicts a disabling condition can be
worth removing early. Concentrating several attacks on its group helps
reduce the number of dangerous actions the party faces in later rounds.

Group attacks are useful when several creatures of the same kind are
present. An all-enemy attack becomes useful when the opposing formation
contains several groups. Check which enemies remain as each new round
begins and choose the next targets accordingly.

\subsection*{Action order and recovery timing}

Agility influences when companions and enemies act. The order can vary
between rounds, so allow for an enemy reaching an injured ally before the
planned heal. A fast companion carrying a Medical Herb can provide another
recovery option when the main healer is slower or needs a different action.

At the start of a difficult round, identify the most urgent job. This may
be restoring a vulnerable ally, curing a disabling condition, reducing
incoming damage, or finishing a dangerous enemy. Assign the remaining
members around that action.

Parry gives a threatened companion a defensive turn while another member
handles recovery. After the party is healthy enough to resume its attack,
return that companion to its usual role. Keep enough recovery available
to handle the next unfavorable sequence of enemy actions.

\subsection*{Choosing to escape}

Run is useful when an encounter threatens the party's remaining supplies
or the journey has already become too expensive in HP and MP. A successful
escape lets the party continue from the field.

A failed attempt gives the enemies an opportunity to act. Consider the
party's condition before attempting another escape, especially when one
more enemy attack could bring down a companion. Healing or protecting the
party can make the next attempt safer.

Some fixed battles and boss encounters restrict escape. Prepare for those
encounters with a recovered party, accessible items, and a clear plan for
the first few rounds.

\section{Physical attacks, critical hits, multi-target weapons, and damage}

\subsection*{Attack power and the enemy's defense}

An ordinary physical attack uses the companion's strength and equipped
weapon against the enemy's physical defense. The resulting damage varies
between strikes. A stronger weapon improves the attacker's opportunities,
and lowering the enemy's defense can make the whole party's physical
attacks more effective.

Watch the damage produced by the current equipment. A sturdy enemy may
be easier to handle after Sap or Decrease has reduced its defense. Bikill
strengthens an ally's physical offense and is useful when that ally will
spend several rounds attacking.

For a short encounter, attacking immediately may finish the group quickly.
In a longer battle, a successful support action has time to improve several
later attacks. Choose the setup according to the opponent's durability
and the party's condition.

\subsection*{Choosing an attack pattern}

Weapon properties determine how far an attack reaches:
\begin{description}
  \item[Single-target weapons:] Ordinary swords, spears, claws, and similar
    weapons concentrate their attack on one creature in the selected group.
    They suit focused attacks against a dangerous enemy or a boss.
  \item[Group weapons:] Whips and weapons such as the Morning Star attack
    grouped enemies. They are useful when several creatures share a target
    entry. The Gringham Whip is a powerful example of this pattern.
  \item[All-enemy weapons:] Boomerangs, including the Bladed Boomerang and
    Flame Boomerang, reach across the enemy formation. The Iron Ball of
    Decimation also attacks all enemies.
  \item[Repeated attacks:] The Falcon Sword and Killer Earring produce
    two attacks from an attack action. Their usefulness depends on the
    wielder's strength, the enemy's defense, and the damage of each strike.
\end{description}

A Boomerang's 15 attack power and a Bladed Boomerang's 25 attack power can
be useful early because their attacks reach the whole opposing formation.
Give a group-clearing weapon to a companion who frequently handles ordinary
encounters. Keep a suitable focused weapon available for a battle with
one durable opponent.

Equipment compatibility remains important. Match the weapon's attack
pattern to a companion who can equip it and use it effectively. After a
change, check the resulting attack value and observe the next few battles.

\subsection*{Critical hits and special weapon effects}

Critical hits produce heavy physical damage and are especially valuable
against enemies with very high defense. They can also make an enemy's
ordinary attack suddenly dangerous. Keep enough HP available to survive
an unfavorable hit when exploring an unfamiliar area.

Some weapons have additional effects. The Poison Needle can produce a
one-hit kill, the Demon Spear also has a possible one-hit-kill effect,
and the Demon Claw can poison an enemy. Their special properties deserve
consideration alongside their attack power and eligible users.

The Dragon Killer is useful against dragons. Other equipment may add an
elemental property or an item-use action. Choose those properties for
the opponents the party is actually fighting, and review the result when
moving to another area.

\section{Spells, special abilities, MP, and reusable battle items}

\subsection*{Matching a spell to the formation}

The Spell command includes the companion's available battle spells and
abilities. Check the target pattern before choosing the action. Blaze
costs 2 MP and targets one enemy, Firebal costs 4 MP and targets a group,
and Bang costs 5 MP and attacks all enemies.

Use a single-target spell to concentrate damage on one important opponent.
A group spell is useful against a cluster of the same species. An all-enemy
spell helps when several groups are present. Enemy resistances affect the
result, so watch how well the selected spell works in the current encounter.

A well-chosen offensive spell can shorten a battle and reduce the healing
needed afterward. Keep the journey's remaining demands in mind when
deciding how much MP to spend on an ordinary group.

\subsection*{Monster abilities}

Monster companions bring abilities such as breath attacks, dances,
physical techniques, and disruption effects. Breath attacks and many
monster-specific actions provide useful options with no MP cost.
Their target patterns and effects still matter when choosing an action.

A breath user can contribute consistent damage through a long dungeon.
A companion with sleep, paralysis, or disruptive abilities can help control
a dangerous group. Inspect the available abilities after recruitment and
after level gains so that the new options become part of the party's plan.

An ability's usefulness depends on the opponent. Switch to another useful
action when a creature resists the chosen effect. Physical damage, another
spell family, support, or an item can give that companion a productive turn.

\subsection*{Healing spells and their costs}

Heal costs 2 MP, Healmore costs 5 MP, and Healall costs 7 MP. Healall restores
one ally's HP completely. Choose a recovery spell that suits the injury
and leaves enough resources for the next rounds.

Healus costs 18 MP and restores HP across the active party. Healusall costs
36 MP and provides full-party HP recovery. These spells become valuable
when an enemy damages several companions at once.

Revival and status treatment add to the healer's workload. Keep another
recovery option available so that the party can handle an injury while
the main healer deals with a fallen or disabled companion.

\subsection*{Reusable item actions}

Certain possessions provide battle effects through the Item command:
\begin{description}
  \item[Staff of Benediction:] Provides a Healmore effect when used in
    battle. It gives its eligible user an additional recovery action.
  \item[Staff of Antimagic:] Provides a Stopspell effect, useful against
    an enemy that depends on spells and is vulnerable to the effect.
  \item[Sage's Stone:] Restores HP to the active party and can be used
    repeatedly during battle.
  \item[Drum of War:] Strengthens the active party's physical offense,
    making it useful when several members will continue attacking.
\end{description}

Give a useful reusable object to an available companion and select its
item action during that companion's turn. Its use occupies an action,
so include that action in the round's plan. A fast user can be useful
when the effect needs to arrive early.

Reusable recovery can preserve a caster's MP for revival, strong healing,
or support. Keep consumables available for urgent situations where another
companion needs to handle the same job.

\section{Status ailments, buffs, debuffs, resistance, and spell reflection}

\subsection*{Responding to disabling effects}

Poison, paralysis, sleep, confusion, and spell sealing can change the
party's priorities in a single round. Check which companions can still
act and which recovery tools are available to them.

Antidote and Antidote items treat poison. Numboff and Full Moon Herbs
treat paralysis. An Angel's Bell provides a way to address confusion.
Carry the relevant items with active members who can use them when an
important caster is disabled.

Sleep and confusion can make a planned action unreliable. Protect the
remaining available companions, keep recovery ready, and remove the
enemy causing the condition when practical. If spell sealing affects
the healer, use accessible item recovery while handling the encounter.

Some enemy effects drain MP. Removing that enemy early can preserve the
resources needed for the rest of the dungeon. Check the healer's remaining
MP after the battle before continuing.

\subsection*{Strengthening the party and weakening an enemy}

Upper raises one ally's physical defense, and Increase strengthens the
party's physical defense. They are useful against an opponent that relies
on ordinary physical attacks. Bikill improves an ally's physical offense.
Choose a companion likely to keep attacking for several rounds when using it.

Sap lowers a single enemy's defense, and Decrease works against a group.
A successful defense reduction can improve the results of several party
members' attacks. A resistant opponent may require a different plan.

Some enemies remove beneficial effects during battle. Watch for that
action and reassess which support effects are worth restoring. Keep
healing and immediate survival ahead of rebuilding a full set of boosts
when the party has just taken heavy damage.

\subsection*{Breath protection and equipment resistance}

Barrier helps protect the active party from breath attacks. It is useful
when a strong fire or cold breath is a regular part of the opponent's
attack pattern. Combine that protection with equipment suited to the
same threat and enough party healing to cover the remaining damage.

Armor and accessories can provide protection from particular attacks or
status effects. Inspect these properties when preparing for an area.
A companion who repeatedly loses its turn to a condition may benefit
from a protective item as well as a ready cure.

Enemy resistance also affects offensive spells and disabling actions.
Watch the battle messages and damage results. A spell that works well
against one species can have a different result against another.

\subsection*{Using Bounce carefully}

Bounce reflects spells and can help protect a companion from a spell-heavy
opponent. Plan recovery for the protected member before relying on it.
Item-based healing provides a useful option while reflection is active.

An enemy using reflection also changes offensive spell selection.
Physical attacks and suitable abilities or item actions give the party
other ways to continue. Check the target's current condition before
committing a powerful spell.

\section{AI tactic modes and when to use direct commands}

\subsection*{Choosing a policy}

Tactics set the behavior of computer-controlled companions. The Hero's
own actions are chosen through the ordinary command menu. Choose the
policy that suits the current opponents, remaining resources, and length
of the journey.

The six policies cover the following approaches:
\begin{description}
  \item[All-out attack:] Prioritizes strong offensive actions and powerful
    available spells or abilities. It suits a group that the party wants
    to remove quickly. Watch MP and the need for a separate healing action.
  \item[Balanced fighting:] Mixes offense with recovery and useful support
    according to the situation. It is convenient for familiar ordinary
    encounters with a healthy, well-equipped party.
  \item[Support the Hero:] Emphasizes assistance for the Hero through
    recovery, strengthening, protection, and hindering the opposition.
    Use it with the Hero active, able to attack, and suitably equipped.
  \item[Conserve MP:] Uses physical attacks and available MP-free actions.
    It is useful during manageable encounters on the way to a demanding
    area. Keep an emergency recovery action available.
  \item[Protect life:] Places greater priority on healing and defensive
    support. It can help a party facing dangerous ordinary groups, and
    its frequent recovery can consume a substantial amount of MP.
  \item[Follow orders:] Gives direct control over eligible companions'
    actions. It is useful for precise targets, planned item use, and
    coordinated boss strategies.
\end{description}

Under an MP-conserving policy, the Hero's manually selected spells still
use their normal MP costs. Check the party's resources as the route
continues and change the policy when its priorities need to change.

\subsection*{When direct orders are useful}

Use Follow orders when a specific companion needs to heal, revive, use
a particular item, or apply a support effect at a planned moment. Direct
orders also help coordinate several attacks against a selected enemy group.

A reusable healing item has a particular owner and uses that owner's turn.
Selecting the action directly makes it easy to account for both. A second
member can handle single-target recovery while the item user restores the
active party.

For a boss encounter, decide the first few rounds in advance. Assign the
initial defense or offense boosts, identify the main recovery member,
and leave room for an urgent cure or revival. Adjust the plan as the
opponent reveals its pattern.

\section{Intelligence, obedience, and monster command restrictions}

\subsection*{Developing reliable companions}

A monster's intelligence affects its ability to follow commands
consistently. A low-intelligence recruit can act on its own or fail to
carry out the intended action. Keep a reliable recovery member available
when training such a companion.

Raising intelligence above 20 provides a practical target for reliable
obedience. Check the Status screen as the monster gains levels and review
its available spells and equipment at the same time.

Intelligence Seeds provide permanent increases. Use them on the intended
companion and check the resulting value after each use. This is useful
for a promising species whose natural intelligence growth or level cap
leaves its command reliability short of what the party needs.

\subsection*{Level caps and a companion's role}

Maximum level varies by species. The Dancing Jewel reaches its cap at
level 7, the Orc King at 30, and the Curer at 50. A Slime Knight can continue
to level 99. Those limits give the companions different growth schedules.

Consider the complete set of abilities available at the current level.
A companion with a low cap can still have useful support or protection,
and a longer-growing species may need several levels before an important
healing spell becomes available.

Keep an unreliable new recruit in a suitable training role while another
member handles essential recovery. Once its intelligence and abilities
support the intended job, move it into a position with more responsibility.

\subsection*{Temporary companions}

Temporary companions such as Papas and Bella have their own battle
behavior. Plan the Hero's actions around the support or offense they
provide. Their autonomous actions remain part of those encounters even
when the ordinary party is using direct orders.

Review the lineup after a temporary companion joins or leaves. Check
the available healing, item owners, and equipment before continuing into
the next difficult area.

\section{The three-member active party and wagon reserves}

\subsection*{Filling the active positions}

The normal active party has three members. With the wagon available,
the travelling roster can contain eight companions, leaving five reserve
positions behind the active three.

Choose a lineup whose abilities work together. Useful jobs include
physical damage, recovery, support, and attacks that reach several enemies.
One companion can cover more than one job, depending on its spells,
equipment, and current level.

For ordinary exploration, a group-clearing attacker, a reliable recovery
member, and a durable companion make a comfortable combination. Before
a stronger opponent, review its likely threats and adjust the roles.
Several physical attackers can benefit from offensive support, while
frequent party-wide damage creates more work for recovery.

\subsection*{Companions with useful combinations}

A Slime Knight combines physical equipment options with healing. It
starts with Heal, learns Healmore at level 5, and gains Healall at level
20. This gives it useful recovery work alongside its attacks.

An Orc King starts with Vivify, learns Healus at level 8, and gains Revive
at level 11. Its developing spell list can cover both revival and
party-wide healing. Keep enough MP available for those jobs when relying
on it during a longer encounter.

A Curer begins with Healall and later learns additional recovery and
support. It gains Increase at level 10, Vivify at 12, Healus at 18, Revive
at 23, and Healusall at 30. Its role can expand substantially as it grows.

Use these combinations when distributing the party's responsibilities.
A companion capable of fighting and healing can cover an urgent injury
while the usual healer revives another member or restores the group.

\subsection*{Making the reserves useful}

Reserve positions can hold another healer, an alternative attacker,
a companion with a particular support ability, or a recruit being trained.
Keep their equipment and condition suitable for the job they may need
to perform.

Where the wagon is available, changing members can bring fresh MP or
a more useful ability into the active lineup. A reserve holding an
important battle item must be brought into an available position to use
it during its turn.

Review both the active party and reserves before leaving town. Healing
and supplies distributed across the travelling roster give more options
during a long sequence of battles.

\section{Changing members, wagon access, experience, and dungeon restrictions}

\subsection*{Changing the lineup}

Use the battle menu's Change command when the wagon and encounter permit
it. Select the available companions needed for the next part of the fight,
then review the new lineup before choosing its actions.

A fresh healer can take over when the original member has spent most of
its MP. A durable attacker can help after a vulnerable companion has
taken heavy damage. A specialist can enter to supply a protective spell
or a particular item action.

The incoming companion needs its own equipment, available MP, and useful
possessions. Prepare these before the fight so that the change provides
an immediate benefit.

\subsection*{Areas with restricted wagon access}

Some areas admit only the active party. Choose the three companions needed
for the dungeon before entering, and give them the supplies they will
need inside. Include reliable healing and a way to handle the area's
important status threats.

Check access when the route moves through an entrance, floor connection,
or fixed encounter. The companions available in an open-field battle can
differ from the members available deeper in a dungeon.

Keep a useful battle item with a companion who will actually enter the
area. The same preparation applies to required keys, recovery supplies,
and any field spell needed for the way out.

\subsection*{Experience and training}

Victories award experience to eligible living companions. In battles
where the wagon is available, reserves can gain experience alongside the
active members. This gives newly recruited companions a way to develop
while a dependable front line handles the encounters.

Choose manageable groups when training a fragile recruit. Keep it alive
and recovered, and inspect its new abilities as levels increase. A useful
healing or support spell may arrive well before its final level cap.

For a dungeon with restricted access, arrange the party around the members
participating in that area. Use accessible wagon battles when the goal
is to develop companions waiting in reserve.

\section{Healing economy, resource conservation, and boss preparation}

\subsection*{Keeping a dungeon run sustainable}

The journey uses HP, MP, and consumables across several encounters. Keep
the main healer's remaining MP visible when deciding how far to explore.
Medical Herbs can handle smaller injuries and preserve spells for the
next difficult stretch.

MP-free abilities and reusable items can provide steady work through a
long route. Assign those actions to companions who can use them reliably,
and keep a stronger recovery option available for a dangerous round.

When escape spells are part of the return plan, reserve their cost.
Outside and Return each cost 8 MP. A companion casting both needs 16 MP
available for that journey home, plus any healing still needed beforehand.

\subsection*{Preparing for a boss}

Rest the party, treat ailments, and check equipment before the approach.
Fill the available inventories with useful recovery and status-treatment
items. Arrange item ownership so that several active members can help
when the regular healer is occupied.

Use the opponent's attack pattern to choose support. Increase helps with
physical pressure, Barrier helps with breath attacks, and suitable armor
can provide further protection. Offensive strengthening becomes useful
when the party can keep an attacker working for several rounds.

Some enemies perform several actions in a round. Allow enough HP for
the possible sequence of attacks before relying on the next heal. Keep
a second recovery action available for a round in which one companion
must be revived or cured.

\subsection*{A workable rhythm for a long fight}

Begin by securing the party's immediate safety and applying useful
support. Continue with focused attacks while watching HP and MP. A
single-target healer can handle an isolated heavy hit, and party-wide
recovery can cover a breath or spell that injures all three members.

The Sage's Stone is useful for repeated party recovery when it is
available. Its user spends an action restoring the group while another
companion handles a more urgent single-target heal, revival, or support
spell.

The Drum of War supports a physically focused active party when available.
After an enemy removes strengthening effects, restore the most useful
ones as the party's condition allows. Keep the recovery plan working
through those changes.

If a healer is running low on MP and the wagon is accessible, bring in
a prepared reserve. Check the incoming member's item ownership and
abilities before relying on the next round.

\section{Enemy summoning, metal enemies, defeat, and recovery}

\subsection*{Handling reinforcements}

Some enemies call additional creatures into the battle. Reinforcements
increase the number of actions the party must handle and can prolong
an encounter that initially looked manageable.

Remove the summoner early when the party can do so safely. Keep track
of the remaining groups and the abilities of newly arrived creatures.
A called healer can make an enemy group harder to finish, while another
disabling enemy can increase the need for recovery items.

Called creatures can also matter during recruitment. Healing monsters
can appear through summoning, so watch for the species being pursued
and manage the battle with that goal in mind. Keep the party healthy
enough to finish the encounter.

\subsection*{Fighting metal enemies}

Metal enemies offer valuable experience and can flee quickly. Their high
defense and strong resistance to offensive magic make physical actions
and critical hits important parts of the plan.

Use direct commands to concentrate suitable attacks on the metal group.
A fast companion has a useful chance to act before the enemy escapes.
Repeated-attack weapons such as the Falcon Sword or Killer Earring can
provide additional physical strikes when an eligible user has them.

Keep ordinary enemies in the encounter under control as well. A harmful
status effect or a badly injured companion can interrupt the attempt.
After a successful battle, check the levels and abilities of companions
who gained experience, including eligible wagon reserves.

Choose a hunting area the party can handle comfortably. Nearby recovery
and dependable equipment help make repeated attempts worthwhile, and
ordinary victories along the route still contribute to development.

\subsection*{Recovering after a difficult encounter}

Treat poison and other persistent conditions, restore enough HP for the
next stretch, and check the remaining MP and items. Reorganize the active
party when another companion is better prepared to continue.

If the party is defeated, half of the carried gold is lost. Banked gold
remains protected. Arrange revival, treatment, and rest after returning
to the recovery point, then reconsider the equipment and supplies for
the next attempt.

Before returning to the difficult area, identify what caused the collapse.
Repeated breath damage may call for protection and group recovery.
Spell sealing makes accessible healing items valuable. A sequence of
heavy physical hits may require better defense, a sturdier companion,
or earlier healing.

Record the recovered adventure at a church once the party is ready.
The next attempt can then begin with its equipment, recovery roles,
and travel plan already in place.

\chapter{First Generation: Childhood}
\label{written-childhood}

The childhood journey begins with Papas close at hand. Townspeople remember
him, sailors speak warmly of him, and his strength makes the roads feel
safe. The Hero's first independent adventures grow out of ordinary village
life: a missing herbalist, a sick friend, and Bianca's determination to help
an animal in trouble.

Take time to explore the houses and speak with the people along the way.
The early finds have immediate value, and the shops offer enough choices
to shape how comfortably the young party handles its first dungeons.
Free resting places make it easy to prepare a little at a time.

\section{Birth scene, the opening ship, and Vista Port}

\subsection*{A name and a voyage}

Choose the Hero's name when starting the adventure. The opening scene
shows Papas waiting anxiously for news of a birth. When he reaches Martha,
the conversation turns to their new son and the name she has chosen.
The scene closes as Martha becomes unwell and Papas calls to her.

The Hero then wakes in a ship's cabin with his father. Papas suggests
some fresh air, and the deck becomes the first place to explore freely.
The crew's conversations establish a little of their history with Papas.
The captain remembers an old kindness, the sailors know that father and
son are travelling together, and the cook has his own opinion about Papas'
cooking.

\subsection*{The cabin's Medical Herb}

Search the dresser in the cabin where the Hero wakes to obtain a
\textbf{Medical Herb}. Keep it for the first independent trip into a cave.
The barrels around the ship are empty, and the chests in the lower crew
quarters are locked during this voyage.

Walk around the deck and speak to the sailors, including the man at the
helm. Visit the captain's cabin and the lower rooms with the cook and
crew. After the conversations, the ship reaches Vista and the captain
asks the Hero to fetch Papas.

Return to the cabin, speak to Papas, and follow him onto the deck.
When he asks whether the drawers have been checked, confirm after collecting
the herb. Father and son disembark, and the ship departs.

\subsection*{Vista Port and the first battle}

At the port, Papas stops to speak with a friend and leaves the Hero free
to wander nearby. The woman upstairs remembers an earlier visit by Papas
and a small child. Explore the building and then leave through the port's
exit when ready to continue.

Three Slimes attack as the Hero steps outside. Choose Attack and let the
encounter play out. Papas arrives during the battle and quickly takes
control of the situation. He heals the Hero afterward and leads the way
to Santa Rosa.

Any encounters on that escorted journey are handled with Papas in the
party. His post-battle healing keeps the Hero ready for the next steps.
Save the Medical Herb for the time when the Hero is exploring alone.

\section{Santa Rosa: Sancho, Bianca, supplies, and the missing herbalist}

\subsection*{A homecoming}

Papas' arrival sets the village in motion. A guard runs ahead with the
news, old acquaintances greet him, and Sancho welcomes both travellers
into his house. Bianca and her mother are visiting while seeking medicine
for Duncan, who is ill at home in Alcapa.

Bianca takes the Hero upstairs and tries to read a book about the sky.
After her mother calls her away, search the room's dresser for a
\textbf{Leather Hat} and equip it on the Hero. This is the first useful
piece of free defensive equipment.

Downstairs, Papas announces that he has business elsewhere. Sancho offers
a place to rest, giving the Hero a convenient recovery point for his
first solo adventure.

\subsection*{Santa Rosa's household finds}

Collect these objects while looking around the village:
\begin{description}
  \item[Leather Hat:] The upstairs dresser in Sancho's house, in the room
    where Bianca reads to the Hero. Equip it immediately.
  \item[Medical Herb:] The dresser on the ground floor of Sancho's house.
  \item[Strength Seed:] The pot in Sancho's basement.
  \item[Medical Herb:] The dresser upstairs in the weapon shop.
\end{description}

The Strength Seed permanently improves strength. The Hero is a useful
recipient because he will be doing the attacking during the cave rescue.
Using it now gives that improvement time to contribute to the early fights.

Together with the ship's herb, the two village herbs provide a small
recovery supply. Sancho's free rest makes it practical to keep those items
for occasions when the Hero is several floors from safety.

\subsection*{Purchase guide: Santa Rosa}

The weapon shop stocks the following early weapons:
\begin{center}
\begin{tabular}{@{}lrr@{}}
\textbf{Weapon} & \textbf{Attack} & \textbf{Price} \\
\hline
Cypress Stick & 2 & 10 G \\
Bamboo Pole & 5 & 50 G \\
Sharpened Bone & 6 & 70 G \\
Oak Cudgel & 9 & 130 G \\
Copper Sword & 13 & 270 G \\
\end{tabular}
\end{center}

A \textbf{Bamboo Pole} is an inexpensive immediate upgrade for 50 G.
It makes the first cave encounters more comfortable while the Hero builds
levels and earns money. The \textbf{Oak Cudgel} is a stronger intermediate
choice if 130 G is already available.

The \textbf{Copper Sword} offers the strongest attack in this shop at
270 G. It is a useful target for players who enjoy a little early training.
Sancho's free rest allows repeated short trips through the upper cave,
so the savings can grow safely. Choose a purchase that fits the current
funds, then equip it before returning to the cave.

The cave itself contains a Leather Shield and Wayfarer's Clothes.
Those finds will improve the Hero's defense during the rescue. Keep the
Leather Hat equipped and focus the village's spending on a useful weapon.

\subsection*{Learning what has happened}

Speak with the villagers and visit Bianca's mother at the inn. The man
who went into the nearby cave to gather medicinal herbs has failed to
return. The household in the mountain north of Sancho's home also has
news of the missing herbalist.

Papas is occupied, and the Hero remains within Santa Rosa at this point.
The cave entrance lies on the west side of town, near the church. Rest
with Sancho, save at the church, and approach the entrance with the new
weapon equipped and the herbs available.

\endgroup

\end{document}